\documentclass[12pt,a4paper]{article}

% ---- Encoding & Font ----
\usepackage{iftex}
\ifPDFTeX
  \usepackage[utf8]{inputenc}
  \usepackage[T1]{fontenc}
\else
  \usepackage{fontspec}  % XeTeX/LuaTeX (tectonic): Unicode nativo (·, —, ¿, ª, §)
\fi
\usepackage[spanish]{babel}
\usepackage{csquotes}

% ---- Page layout ----
\usepackage[top=2.5cm, bottom=2.5cm, left=2.5cm, right=2.5cm]{geometry}
\usepackage{setspace}
\onehalfspacing

% ---- Math ----
\usepackage{amsmath, amssymb, amsthm}
\usepackage{mathtools}
\usepackage{physics}

% ---- Graphics ----
\usepackage{graphicx}
\usepackage{tikz}
\usepackage{float}
\usepackage{caption}

% ---- Tables ----
\usepackage{array}
\usepackage{booktabs}
\usepackage{tabularx}
\usepackage{colortbl}

% ---- Colours ----
\usepackage{xcolor}
\definecolor{notebg}{HTML}{FFF3CD}
\definecolor{noteborder}{HTML}{FFC107}
\definecolor{boxred}{HTML}{F8D7DA}
\definecolor{boxredborder}{HTML}{F5C6CB}

% ---- Boxes ----
\usepackage{mdframed}
\usepackage{tcolorbox}
\tcbuselibrary{most}

\newtcolorbox{keybox}[1][]{
  colback=blue!5,
  colframe=blue!70!black,
  arc=4pt,
  boxrule=1pt,
  fonttitle=\bfseries,
  title={#1}
}

\newtcolorbox{warnbox}[1][]{
  colback=boxred,
  colframe=boxredborder,
  coltitle=black!75,
  arc=4pt,
  boxrule=1pt,
  fonttitle=\bfseries,
  title={#1}
}

\newtcolorbox{notebox}[1][]{
  colback=notebg,
  colframe=noteborder,
  coltitle=black!75,
  arc=4pt,
  boxrule=1pt,
  fonttitle=\bfseries,
  title={#1}
}

% ---- Hyperlinks & TOC ----
\usepackage[hidelinks]{hyperref}
\usepackage{bookmark}

% ---- Enumerate ----
\usepackage{enumitem}

% ---- Miscellaneous ----
\usepackage{icomma}
\usepackage{siunitx}
\usepackage{textcomp}
\usepackage[bottom]{footmisc}

% ---- Title ----
\title{\textbf{Clase 4: Los números reales} \\[0.3em]
        \large{Axiomática de cuerpo ordenado: suma, producto y orden}}
\author{Transcripto y expandido --- Curso Cálculo 1 (OpenFING)}
\date{}

\begin{document}

\maketitle
\thispagestyle{empty}
\newpage

\tableofcontents
\newpage

\section{Introducción: la recta real como intuición geométrica}

La clase abre con un cambio de punto de vista respecto al liceo: en vez de
dar por sentado qué es $\mathbb{R}$, se pregunta por qué es un tema difícil.
La intuición fundamental es que los \textbf{números reales son los números
que permiten representar todos los puntos de una recta} --- y la palabra
clave es \emph{todos}. Fijado un sistema de coordenadas en una recta (un
punto $O$ para el cero y un punto $I$ para la unidad), cada número real $x$
se asocia biunívocamente a un punto $X$ tal que $\vec{OX} = x \cdot
\vec{OI}$.

\begin{notebox}[Biyección con toda recta]
Esta correspondencia es \textbf{biyectiva}: hay tantos reales como puntos en
la recta. En consecuencia, \textbf{todas las rectas están en biyección con
$\mathbb{R}$}: todas tienen la misma cantidad de puntos. Por eso se habla de
\emph{la} recta real, entendida como una recta ideal que representa a todas
las rectas posibles.
\end{notebox}

$\mathbb{R}$ contiene, por supuesto, a los naturales, los enteros y los
racionales. Pero ya se sabe que hay más puntos en la recta que racionales:
el ejemplo clásico es $\sqrt{2}$, construible con compás como la diagonal de
un cuadrado de lado $1$, y que no es racional (demostración diferida a una
clase futura). A esta lista se agregan $\pi$, $e$, $\log 2$, etc. Aunque
parezca paradójico escribir solo cuatro irracionales frente a infinitos
racionales, es un resultado de la segunda mitad del siglo XIX que
\textbf{el infinito de los irracionales es estrictamente más grande} que el
de los racionales (la herramienta para probarlo --- cardinalidad --- no está
disponible todavía en el curso).

\begin{notebox}[Nota histórica: griegos antiguos]
Los griegos antiguos distinguían dos tipos de objetos: los \emph{números}
(solo los racionales) y los \emph{objetos geométricos} (que cubrían todos
los reales). Como no tenían notación para los irracionales, Euclides
representa siempre los números por puntos de recta --- de ahí que los
\emph{Elementos} estén llenos de figuras.
\end{notebox}

\section{Dos caminos hacia una definición rigurosa: construcción y axiomatización}

El cálculo diferencial e integral nace a fines del siglo XVII con
\textbf{Leibniz}, quien inventó de una vez el cálculo diferencial, el
integral, y buena parte de lo que se ve en Cálculo 1 y 2 --- sin tener, sin
embargo, una definición precisa de número real, ni noción de límite, ínfimo
o supremo. Esa carencia fue el motor que llevó, en el siglo XIX, a precisar
la noción de número real por dos vías:

\begin{itemize}
  \item \textbf{Construcción}: partir de ``la nada'' (o de los racionales
    $\mathbb{Q}$) y construir un objeto que funcione como $\mathbb{R}$. El
    paso de los enteros a $\mathbb{Q}$ es sencillo; el paso difícil de
    $\mathbb{Q}$ a $\mathbb{R}$ lo resolvió el matemático alemán
    \textbf{Dedekind} (cortes de racionales), y también hubo una
    construcción alternativa de \textbf{Cauchy} (sucesiones).
  \item \textbf{Axiomatización}: capturar las propiedades de $\mathbb{R}$
    mediante una lista de fórmulas sencillas, al estilo de los axiomas de la
    geometría de Euclides (siglo V-VI a.C.).
\end{itemize}

Lo notable es que \textbf{ambos programas confluyeron}: se demostró que las
construcciones de Dedekind y Cauchy satisfacen todos los axiomas de la
axiomatización, y que \textbf{si existe un objeto que satisface la lista de
axiomas, es único a menos de isomorfismo}. La analogía con informática que
da el docente: así como una misma estructura de datos (un conjunto finito,
por ejemplo) puede implementarse con listas o con árboles binarios sin que
cambien sus propiedades observables, todas las construcciones posibles de
$\mathbb{R}$ son isomorfas entre sí.

\begin{notebox}[Por qué se admite $\mathbb{R}$ sin construirlo]
Por esta razón el curso admite la existencia de un único cuerpo
$\mathbb{R}$ que satisface la lista de axiomas, sin presentar la
construcción (``es gigantesca'').
\end{notebox}

\section{La estructura de $\mathbb{R}$: quince axiomas en cuatro grupos}

Se postula que existe un conjunto $\mathbb{R}$ dotado de dos
\textbf{operaciones}, $x+y$ (suma) y $x\cdot y$ (producto, también escrito
$xy$), y una \textbf{relación} binaria $x\le y$ (equivalente a escribir
$y\ge x$ --- mera notación sinónima, no una definición nueva), que satisface
\textbf{15 axiomas}, agrupados así:

\begin{table}[H]
\centering
\begin{tabular}{@{}clp{7cm}@{}}
\toprule
Grupo & Rango & Tema \\
\midrule
1 & 1.1 -- 1.4 & Suma \\
2 & 2.1 -- 2.5 & Producto \\
3 & 3.1 -- 3.5 & Orden \\
4 & único & Completitud (motor de toda la teoría) \\
\bottomrule
\end{tabular}
\end{table}

\subsection{Axiomas de la suma}

\begin{align}
\textbf{1.1 (asociatividad)} &\quad (x+y)+z = x+(y+z), \ \forall x,y,z \\
\textbf{1.2 (conmutatividad)} &\quad x+y = y+x, \ \forall x,y \\
\textbf{1.3 (neutro)} &\quad \exists\, 0 : \ x + 0 = x, \ \forall x \\
\textbf{1.4 (opuesto)} &\quad \forall x, \ \exists\, (-x) : \ x + (-x) = 0
\end{align}

\subsection{Axiomas del producto}

\begin{align}
\textbf{2.1 (asociatividad)} &\quad (xy)z = x(yz) \\
\textbf{2.2 (conmutatividad)} &\quad xy = yx \\
\textbf{2.3 (neutro)} &\quad \exists\, 1 \ne 0 : \ x \cdot 1 = x, \ \forall x \\
\textbf{2.4 (inverso)} &\quad \forall x \ne 0, \ \exists\, x^{-1} : \ x\cdot x^{-1} = 1 \\
\textbf{2.5 (distributividad)} &\quad (x+y)z = xz + yz
\end{align}

\begin{warnbox}[El axioma 2.3 dice más de lo que parece]
No solo exige que exista un neutro para el producto, sino que sea
\textbf{distinto de $0$}. Sin esa cláusula sería trivial demostrar que
$\mathbb{R}$ tiene un único punto (todo sería igual a $0$ y todas las demás
reglas seguirían cumpliéndose vacuamente). Es, en la práctica, el axioma que
garantiza que $\mathbb{R}$ sea infinito.
\end{warnbox}

\subsection{Axiomas de orden}

\begin{align}
\textbf{3.1 (total)} &\quad x\le y \ \text{ó} \ y \le x, \ \forall x,y \\
\textbf{3.2 (transitivo)} &\quad x\le y \ \text{y} \ y\le z \implies x\le z \\
\textbf{3.3 (antisimétrico)} &\quad x\le y \ \text{y} \ y\le x \implies x=y \\
\textbf{3.4 (comp.\ con suma)} &\quad x\le y \implies x+z \le y+z \\
\textbf{3.5 (comp.\ con producto)} &\quad x\ge 0 \ \text{y} \ y\ge 0 \implies xy \ge 0
\end{align}

\begin{notebox}[El ``ó'' es inclusivo]
En 3.1, el ``ó'' es inclusivo: cuando $x=y$ ambas desigualdades se cumplen a
la vez, y eso no es una contradicción.
\end{notebox}

\subsection{El axioma que falta: completitud}

Queda un axioma más, el \textbf{4}, el axioma de \textbf{completitud} ---
descrito como ``el motor de la máquina''. No se presenta en esta clase; se
pospone a la próxima semana. Sin él se puede razonar con las reglas
algebraicas y de orden usuales, pero \textbf{no} se puede todavía justificar
la existencia de $\sqrt{2}$ ni de ningún ínfimo o supremo: esa es
precisamente la razón de ser del axioma 4.

\section{Consecuencias de los axiomas de la suma}

\subsection{Unicidad del neutro y del opuesto}

El axioma 1.3 dice que \emph{existe} un neutro, pero no dice que sea único:
hay que demostrarlo. Supongamos que $0$ y $0'$ son ambos neutros para la
suma. Entonces
\[
0 + 0' = 0' \quad (\text{porque } 0 \text{ es neutro}), \qquad
0 + 0' = 0 \quad (\text{porque } 0' \text{ es neutro}),
\]
de donde $0=0'$. Por esto se puede hablar \emph{del} neutro $0$ (una
constante), y no de \emph{un} neutro cualquiera. El mismo argumento
(dejado como ejercicio) da la \textbf{unicidad del opuesto}.

De la unicidad del opuesto se obtiene la siguiente proposición, usada
constantemente en lo que sigue:

\begin{keybox}[Caracterización del opuesto]
\[
\boxed{x+y = 0 \iff x = -y \iff y = -x}
\]
\end{keybox}

La primera equivalencia y la segunda son la misma afirmación leída con los
papeles de $x$ e $y$ intercambiados (por conmutatividad); solo hace falta
demostrar una de las dos implicaciones no triviales, lo cual queda como
ejercicio. De esta proposición se deducen inmediatamente tres consecuencias
muy usadas:
\[
-(x+y) = (-x) + (-y), \qquad -(-x) = x, \qquad -0 = 0.
\]

\subsection{La resta}

La resta \textbf{no es un axioma}: se \textbf{define} a partir de la suma y
el opuesto.

\begin{keybox}[Definición de la resta]
\[
\boxed{x - y \ := \ x + (-y)}
\]
\end{keybox}

Todas las propiedades usuales de la resta (las que se manejan desde la
primaria) se deducen de esta definición combinada con las propiedades del
opuesto y de la suma --- el ejercicio propuesto es tomar la lista de reglas
algebraicas conocidas y deducirlas una por una a partir de los axiomas.

\section{Consecuencias de los axiomas del producto}

\subsection{Unicidad del neutro y del inverso}

El mismo argumento de la sección anterior (sustituyendo $+$ por $\cdot$ y
$0$ por $1$) demuestra que el neutro $1$ del producto es único, y de la
misma manera el inverso $x^{-1}$ de un elemento no nulo es único.

\subsection{Cero es absorbente: la demostración completa}

Ningún axioma dice explícitamente que $0\cdot x = 0$: \textbf{es una
consecuencia}, y la clase la demuestra paso a paso como ejemplo de cómo se
``juega'' con la lista de axiomas (``un juego infantil... pero solo usa los
axiomas''):
\begin{align}
0x &= 0 + 0x &&\text{(1.3: } 0 \text{ es neutro, sumado a la izquierda)} \\
   &= (0+0)x &&\text{(1.3: } 0+0=0\text{, sustituyendo)} \\
   &= 0x + 0x &&\text{(2.5: distributividad)}
\end{align}
Combinando las dos cadenas se obtiene $0x = 0x+0x$. Sumando $-(0x)$ a ambos
lados y usando 1.4,
\[
0 = 0x + 0x - 0x = 0x.
\]

\begin{keybox}[Cero es absorbente]
\[
\boxed{x \cdot 0 = 0 \cdot x = 0}
\]
\end{keybox}

\begin{notebox}[Modelo de demostración]
Este es el primer ejemplo completo de demostración a partir de la lista de
axiomas, y sirve de modelo para todas las demás: nunca se usa nada que no
esté en la lista.
\end{notebox}

\subsection{Producto nulo y producto unitario}

A partir de la absorción de $0$ se puede demostrar (queda para ejercicio) la
propiedad central:

\begin{keybox}[Producto nulo]
\[
\boxed{xy = 0 \iff x=0 \ \text{ó} \ y=0}
\]
\end{keybox}

y, de modo simétrico a la proposición para la suma,
\[
xy = 1 \iff x = y^{-1} \iff y = x^{-1}.
\]
Como $1\ne 0$ (axioma 2.3), si $xy=1$ entonces ni $x$ ni $y$ pueden ser
nulos, así que ambos son invertibles. De aquí se deducen todas las reglas
usuales de manipulación de productos e inversos.

\section{Los naturales como abreviaturas}

Al desarrollar $(x+y)^2$ con la distributividad aparece el símbolo $2$, que
\textbf{no está definido por ningún axioma}. Se introduce entonces por
definición:

\begin{keybox}[Definición de los enteros mayores que 1]
\[
\boxed{2 := 1+1, \quad 3 := 2+1, \quad 4:=3+1, \ \ldots}
\]
\end{keybox}

Cada entero mayor que $1$ es, formalmente, una \textbf{abreviatura}
construida a partir del $1$ del axioma 2.3 y la suma. El docente ilustra que
estas abreviaturas son consistentes con las reglas que ya se conocen
probando que $x+x=2x$:
\begin{align}
x + x &= 1\cdot x + 1\cdot x && \text{(2.3: neutro)} \\
&= (1+1)x && \text{(2.5: distributividad)} \\
&= 2x && \text{(definición de } 2\text{)}
\end{align}

Con esta maquinaria, el desarrollo del cuadrado de un binomio queda
completamente justificado:
\[
(x+y)^2 = (x+y)(x+y) = x^2 + xy + yx + y^2 = x^2 + 2xy + y^2.
\]

\section{El cociente y la reducción al mismo denominador}

\subsection{Definición y el peligro de dividir entre cero}

Al igual que la resta, el cociente \textbf{se define}, no es un axioma:

\begin{keybox}[Definición del cociente]
\[
\boxed{\dfrac{x}{y} := x \cdot y^{-1}, \quad \text{si } y \ne 0}
\]
\end{keybox}

\begin{warnbox}[No dividir sin verificar el denominador]
Escribir $y^{-1}$ (o $x/y$) sin verificar antes que $y\ne 0$ es un error
grave --- ``se van directamente al infierno... antes de dar el examen'',
con el castigo cómico de repetir el curso de Cálculo 1 por la eternidad. El
chiste subraya un punto serio: $0$ no tiene inverso, y toda manipulación
con cocientes debe justificar explícitamente que el denominador no se
anula.
\end{warnbox}

\subsection{Ejemplos: racionales e irracionales}

De la definición de cociente y de las propiedades del producto se deduce la
regla de \textbf{reducción al mismo denominador}, familiar desde la
primaria:

\begin{keybox}[Reducción al mismo denominador]
\[
\boxed{\dfrac{x}{a} + \dfrac{y}{b} = \dfrac{xb + ya}{ab}}
\]
\end{keybox}

El docente insiste en un punto conceptual importante: esta regla \textbf{no
es exclusiva de los racionales}. Se cumple igual para cualquier par de
reales, con denominadores racionales o irracionales.

\textbf{Caso racional} (el ``ejemplo idiota''):
\[
\frac{1}{2}+\frac{1}{3} = \frac{3}{6}+\frac{2}{6} = \frac{5}{6}.
\]

\textbf{Caso irracional}, anticipando la existencia de $\sqrt2$ y $\sqrt3$
(garantizada en rigor solo por el axioma de completitud, todavía no
presentado, pero ya utilizable informalmente):
\[
\frac{1}{\sqrt2}+\frac{1}{\sqrt3} = \frac{\sqrt3}{\sqrt2\sqrt3}+\frac{\sqrt2}{\sqrt2\sqrt3} = \frac{\sqrt2+\sqrt3}{\sqrt6}.
\]
Para eliminar la raíz del denominador, se multiplica numerador y
denominador por $\sqrt6$:
\[
\frac{\sqrt2+\sqrt3}{\sqrt6} = \frac{(\sqrt2+\sqrt3)\sqrt6}{6} = \frac{\sqrt{12}+\sqrt{18}}{6}
= \frac{2\sqrt3 + 3\sqrt2}{6},
\]
usando $\sqrt{12}=\sqrt{4\cdot3}=2\sqrt3$ y $\sqrt{18}=\sqrt{9\cdot2}=3\sqrt2$.

\begin{notebox}[Idea a retener]
La reducción al mismo denominador es una consecuencia algebraica pura de
los axiomas del producto y de la definición de cociente; por eso funciona
igual de bien racionalizando denominadores irracionales.
\end{notebox}

\section{Interpretación geométrica de las operaciones}

Antes de pasar a los axiomas de orden, se hace una pausa para mostrar que
\textbf{toda cuenta algebraica tiene un correlato geométrico} --- el
mensaje explícito es que quien maneja los números reales solo como
``fórmulas esotéricas'' o ``magia negra'' se pierde, mientras que quien los
piensa como puntos de una recta tiene intuición (``magia blanca'').

\subsection{La suma: construcción por paralelogramos}

Dados $0$, $x$ y $y$ en la recta $R$, se construye el punto suma $z=x+y$ sin
necesidad de conocer la unidad (la suma es una operación \textbf{lineal},
no requiere escala):

\begin{enumerate}
  \item Se toma una recta paralela $R'$ y un punto $O'$ arbitrario en ella.
  \item Se copia el segmento de referencia formando el paralelogramo
    $O'I'IO$, lo que define $I'$.
  \item Se repite la construcción de paralelogramos para trasladar el
    punto $y$ a partir de $x$, obteniendo $z$.
\end{enumerate}

Es la construcción a regla y compás de la suma mediante traslaciones de
vectores (paralelogramos). Queda como ejercicio la construcción análoga
para el opuesto.

\subsection{El producto: el teorema de Tales}

A diferencia de la suma, el producto \textbf{sí necesita la unidad} $I$.
Dados $O$, $I$, $x$ e $y$ sobre la recta $R$:

\begin{enumerate}
  \item Se traza una recta secante $R'$ por $O$ y se marca en ella un
    punto arbitrario $I'$ (la ``unidad duplicada'').
  \item Se traza el segmento $I'I$ y se ubica, sobre $R'$, un punto llamado
    $G'$ trazando la paralela adecuada, de modo que por el
    \textbf{teorema de Tales} se establece una razón de referencia entre
    los segmentos $OI'$, $OG'$ y $OI$.
  \item Repitiendo la construcción entre $I'$ y $X$ (trazando la paralela
    correspondiente) se obtiene un punto $Z$ tal que, otra vez por Tales,
    \[
    \frac{OI'}{OG'} = \frac{OX}{OZ},
    \]
    y $Z$ es exactamente el punto que representa $x\cdot y$.
\end{enumerate}

\begin{notebox}[Paralelogramos vs.\ triángulos semejantes]
La diferencia conceptual entre las dos construcciones: la suma usa
\textbf{paralelogramos} (geometría afín, sin necesidad de escala); el
producto usa \textbf{triángulos semejantes y el teorema de Tales}
(geometría con escala, por eso hace falta la unidad).
\end{notebox}

Como ejercicio se propone construir, con la misma idea, el \textbf{inverso}
$y^{-1}$ a partir de $O$, $I$ e $y$: se trata de invertir los papeles en la
relación de Tales de modo que el producto de $y$ por el punto buscado dé
$1$.

\subsection{Adelanto: la raíz cuadrada con regla y compás}

Como anticipo (no demostrado con los axiomas vistos, solo ilustrado
geométricamente): dados $O$ (cero), $I$ (uno) y $X$ sobre la recta, se marca
el simétrico $J$ de $X$ respecto de $O$ (que representa $-1$ en la escala de
$X$). Tomando el segmento $JX$ como \textbf{diámetro de una circunferencia}
y trazando la perpendicular a la recta por $O$, el punto de intersección
$Z$ con la circunferencia cumple
\[
OZ = \sqrt{OX},
\]
una construcción clásica basada en el teorema del cateto / ángulo inscrito
en semicircunferencia.

\begin{warnbox}[Esta construcción todavía no está justificada]
La construcción de la raíz cuadrada \textbf{no se puede justificar todavía}
con los 15 axiomas vistos, porque la existencia de la raíz cuadrada es
consecuencia del axioma de completitud (axioma 4), que se verá recién la
próxima clase. Se muestra únicamente para ilustrar el punto de vista griego
antiguo: sin notación algebraica para los irracionales, Euclides (por
ejemplo, en la demostración de que hay infinitos primos) razonaba
enteramente con figuras.
\end{warnbox}

\section{Axiomas de orden en detalle}

\subsection{Totalidad, transitividad, antisimetría}

\begin{itemize}
  \item \textbf{3.1 (total)}: dados $x,y$ cualesquiera, siempre se puede
    comparar: $x\le y$ o $y\le x$. Tomando $x=y$ se obtiene como caso
    particular la \textbf{reflexividad}: $x\le x$ para todo $x$.
  \item \textbf{3.2 (transitivo)}: $x\le y$ y $y\le z \implies x\le z$.
  \item \textbf{3.3 (antisimétrico)}: $x\le y$ y $y\le x \implies x=y$.
\end{itemize}

\begin{notebox}[El orden no siempre es total: la inclusión]
No todo orden es total: la \textbf{inclusión entre conjuntos} es un orden
(transitivo, reflexivo, antisimétrico) que \textbf{no} es total --- dados
dos conjuntos cualesquiera, no siempre uno está incluido en el otro. En
cambio, el orden de $\mathbb{R}$ sí es total, precisamente porque
corresponde a la disposición de puntos sobre una recta (``a la izquierda
de'' siempre se puede decidir).
\end{notebox}

\begin{keybox}[El axioma más importante: antisimetría como método de prueba]
El axioma 3.3 es, según el docente, ``quizás el axioma de orden más
importante de todo el universo'': parte de hipótesis sobre el
\textbf{orden} y concluye una \textbf{igualdad}. Gracias a él aparece un
\textbf{método nuevo para demostrar igualdades} que no depende de
manipulaciones algebraicas: para probar $x=y$ basta mostrar por separado
$x\le y$ \textbf{y} $y\le x$. Este recurso se usará con mucha frecuencia
más adelante en el curso.
\end{keybox}

\subsection{Orden estricto y tricotomía}

El orden estricto se define a partir del amplio (y viceversa):

\begin{keybox}[Orden estricto]
\[
\boxed{x < y \ :\Longleftrightarrow\ x\le y \ \text{y} \ x\ne y}
\qquad\qquad
x \le y \iff (x<y \ \text{ó} \ x=y).
\]
\end{keybox}

De aquí se deducen las propiedades del orden estricto:

\begin{itemize}
  \item \textbf{Irreflexividad}: nunca se cumple $x<x$.
  \item \textbf{Transitividad}: $x<y$ y $y<z \implies x<z$.
  \item \textbf{Tricotomía}: para todo par $x,y$, se cumple
    \textbf{exactamente una} de las tres opciones
    \[
    x<y \quad \text{ó} \quad x=y \quad \text{ó} \quad x>y,
    \]
    y aquí el ``ó'' es, por primera vez en la clase, \textbf{exclusivo}.
\end{itemize}

La tricotomía permite clasificar a $\mathbb{R}$ en tres partes disjuntas:

\begin{keybox}[Clasificación por signo]
\[
\boxed{\text{positivos } (x>0), \quad \text{negativos } (x<0), \quad \text{nulo } (x=0)}
\]
\end{keybox}

\section{Un aparte de lógica: la implicación material}

Al discutir por qué la transitividad ($x<y \wedge y<z \implies x<z$)
\textbf{no} es un ``si y solo si'', surge la pregunta de por qué la
implicación no se puede invertir. El docente construye un contraejemplo con
valores concretos:
\[
\underbrace{1 < -3}_{\text{falso}} \quad \text{y} \quad \underbrace{-3 < 2}_{\text{verdadero}}
\qquad \Longrightarrow \qquad \underbrace{1<2}_{\text{verdadero}}.
\]

La hipótesis (falso $\wedge$ verdadero $=$ falso) implica la conclusión
(verdadero), y la implicación completa es \textbf{verdadera} --- porque
\textbf{lo falso implica cualquier cosa}. Esto no es un capricho: es la
definición estándar de la implicación material, y sirve para entender por
qué un enunciado universal ($\forall x,y,z$) puede evaluarse como
verdadero aun sustituyendo valores que hacen falsa la hipótesis.

\begin{notebox}[Puntos remarcados por el docente]
\begin{itemize}
  \item La matemática \textbf{no prohíbe lo falso}: una proposición puede
    ser verdadera aunque contenga subfórmulas falsas, siempre que la
    implicación global sea verdadera.
  \item Si a la izquierda de un ``$\implies$'' hay algo falso, a la derecha
    puede haber \textbf{cualquier cosa} (verdadero o falso) sin violar la
    implicación.
  \item Lo verdadero, en cambio, \textbf{nunca puede implicar} lo falso ---
    por eso la recíproca de una implicación válida no tiene por qué
    cumplirse.
\end{itemize}
\end{notebox}

\section{Orden compatible con la suma y con el producto}

\subsection{Suma: sin restricciones}

Del axioma 3.4 (y su combinación con los demás) se obtiene una equivalencia
completa, sin ninguna condición sobre $z$:

\begin{keybox}[Compatibilidad del orden con la suma]
\[
\boxed{x\le y \iff x+z \le y+z \iff x-z \le y-z, \quad \text{para todo } z}
\]
\end{keybox}

y, al tomar el opuesto, el orden se invierte:
\[
x\le y \iff -x \ge -y.
\]
Las mismas equivalencias valen con desigualdad estricta. El docente lo
resume como ``la fiesta'': sumar o restar una cantidad cualquiera a ambos
lados de una desigualdad \textbf{nunca} requiere verificar nada sobre esa
cantidad.

\subsection{Producto: el origen del 90\% de los errores}

Con el producto la historia es distinta, y el docente es explícito: es
\textbf{responsable del 90\% de los errores de los estudiantes} en
Cálculo 1. La tabla resume los casos, todos derivados del axioma 3.5 y de
la tricotomía:

\begin{table}[H]
\centering
\begin{tabular}{@{}lp{8cm}@{}}
\toprule
Hipótesis sobre $z$ & Conclusión a partir de $x\le y$ \\
\midrule
$z \ge 0$ & $xz \le yz$ (se conserva el sentido) \\
$z \le 0$ & $xz \ge yz$ (se \textbf{invierte} el sentido) \\
$z = 0$ & $xz = yz = 0$ (caso degenerado, ninguna desigualdad estricta sobreviene) \\
\bottomrule
\end{tabular}
\end{table}

Con desigualdad \textbf{estricta} el caso $z=0$ deja de ser inofensivo: si
$x<y$ y $z=0$, entonces $xz=yz$ (una \textbf{igualdad}, no una desigualdad
estricta), así que la implicación estricta exige $z>0$ (estrictamente):

\begin{keybox}[Compatibilidad del orden con el producto (versión estricta)]
\[
\boxed{x<y \ \text{y} \ z>0 \implies xz<yz}
\]
\[
x<y \ \text{y} \ z<0 \implies xz>yz.
\]
\end{keybox}

\begin{warnbox}[Advertencia central de la clase]
Nunca se debe multiplicar (ni dividir) una desigualdad por una cantidad sin
\textbf{antes} declarar y justificar su signo. El docente insiste en que
buena parte de las pruebas del primer parcial contendrá pasos de este tipo
hechos sin esa verificación, y que es exactamente el tipo de error que hace
perder el curso.
\end{warnbox}

\emph{Próxima clase: el axioma 4 (completitud), el que permite finalmente
justificar la existencia de $\sqrt{2}$, del ínfimo y del supremo.}

\end{document}
